\subsection{A coherent Gaussian convolution carrying the actual tensor maps}
Finding~\fid{F50} supplied exact Gaussian path signs and a faithful
realization of point-Weil Hom spaces. Its failure of ordinary tensor
monoidality has an explicit geometric repair. The repair uses a
specified convolution and leaves the pointwise weight question unchanged.

\begin{finding}[\fid{F53}: geometric convolution and actual tensor operations]
\label{f:gaussianconvolution}
Retain the odd-characteristic Gaussian data of \fid{F50}:
\[
 L=\mathcal L_\psi(x^2),\qquad
 G=H_c^1(\A^1_{\bar k},L),\qquad
 J(C)=L\otimes\underline{C\otimes G^\vee}[1].
\]
An underline denotes constant pullback from the finite-field point.
There is an actual geometric product
\begin{equation}\label{eq:gaussianstar}
 A\star_\psi B=
 \underline{G(1)}[1]\otimes
 Rm_!\left(A\boxtimes B\otimes\mathcal L_\psi(2xy)\right),
 \qquad m(x,y)=x+y,
\end{equation}
with the associator and symmetry specified below, for which
\[
 J(C)\star_\psi J(C')\simeq J(C\otimes C').
\]
On the full perverse image of $J$, its unit is $J(E)$ and $J$ is
an exact fully faithful symmetric strong monoidal equivalence from
finite Weil spaces. Its image unit differs from its ambient unit,
and its transported symmetry differs by a sign from the naive
kernel swap on two Gaussian image objects.

Composing $J$ with \fid{F48}'s actual fixed-parameter tensor
functor gives an exact faithful strong monoidal realization of the
nonstandard family through all degrees by pure weight-zero
perverse objects. It preserves the actual named maps and their
existing relations. The original external-product path construction
of \fid{F50} retains its exact negative effective spectrum and is
natural in those coefficient maps.

This supplies coherent operations on these chosen objects.
It does not make the source functor full, identify its quantum
symmetry with the new symmetry, select the designated basis, or
improve the exact pointwise/dual estimate of \fid{F50}.
\end{finding}

\begin{proof}
\par\smallskip\noindent\textbf{Normalized constants.}\quad
Let $*$ be ordinary additive convolution
$A*B=Rm_!(A\boxtimes B)$ on constructible Weil complexes.
It has the usual Fubini associator, the derived symmetry
induced by variable interchange, and ambient unit $\delta_0$.
Put
\[
 \Theta=\underline E(1)[2],\qquad
 K(C)=\underline C\otimes\Theta.
\]
The coordinates $(s=x+y,y)$ make $m$ a trivial affine-line bundle.
Consequently
\[
 Rm_!\underline E=\underline E(-1)[-2].
\]
Canonical compact trace gives natural isomorphisms
\begin{equation}\label{eq:normalizedconstantconv}
 K(C)*K(C')\simeq K(C\otimes C').
\end{equation}
Explicitly, its left side is
$\underline{C\otimes C'}\otimes Rm_!\underline E(2)[4]$.

The primary formalism is affine-space cohomology
\cite[Example~16.3, p.~109]{Milne},
base change and compactification
\cite[Theorem~17.10, p.~116; Proposition~18.4,
pp.~121--122]{Milne}, and compact Kunneth and projection
formula \cite[Theorem~22.1, Remarks~22.3(b),(d),
Proposition~22.5, pp.~135--137]{Milne}.
The canonical trace
\[
 R\Gamma_c(\A^1,\underline E(1)[2])\simeq E
\]
is normalized by sending the class of each geometric point
to one \cite[Theorem~24.1(a),(b), p.~144]{Milne}.

We check coherence directly. For the addition map $m_r$
on $r$ coordinates, the fibers are affine $(r-1)$-spaces.
Their canonical relative compact trace gives
\begin{equation}\label{eq:normalizedconstantmulti}
 R(m_r)_!\bigl(K(C_1)\boxtimes\cdots\boxtimes K(C_r)\bigr)
 \simeq K(C_1\otimes\cdots\otimes C_r).
\end{equation}
Every iterated binary comparison is this map.
Indeed the fiber has compact cohomology only in top degree
$2(r-1)$, successive normalized traces send its point
class to one, and Kunneth sends a product of point classes
to the point class in the product.
Changes between fiber coordinates for different parenthesizations
are affine automorphisms. They preserve the canonical trace,
since they send point classes to point classes.
Thus every parenthesization gives
\eqref{eq:normalizedconstantmulti}; this proves the pentagon
and all higher tensor compatibilities.

Taking either coefficient to be $E$ in
\eqref{eq:normalizedconstantconv} gives left and right unit
constraints with unit $\Theta$. The three-factor comparison
proves the triangle identity.
There is also a concrete ambient unit map.
For $i:\{0\}\hookrightarrow\A^1$, cohomological purity and
the $i_*$--$Ri^!$ adjunction give
\[
 Ri^!E=E(-1)[-2],\qquad \delta_0\longrightarrow\Theta.
\]
This is the origin Gysin class
\cite[Theorem~16.1, p.~108; Theorem~16.7 and its
preceding adjunction, p.~111]{Milne}.
Convolution with this map on a normalized constant object
is inverse to its unit comparison: the composite integrates
a point class and is therefore the identity.

Permutations of the inputs preserve the affine-fiber trace.
The shifts $[2]$ are even, so the resulting coefficient
symmetry is the ordinary flip with no additional sign.
The same $r$-factor comparison proves its compatibility
with the associator.
Equal shifts and twists cancel in degree-zero Hom;
connectedness of geometric $\A^1$ identifies constant-sheaf
maps with linear maps, and Weil compatibility is Frobenius
commutation.
Hence $K$ is an exact fully faithful symmetric strong
monoidal equivalence onto the full normalized-constant
subcategory, with its unit $\Theta$.
This unit has full support and is not the ambient unit
$\delta_0$.

\par\smallskip\noindent\textbf{The Gaussian gauge and both units.}\quad
Define an autoequivalence
\[
 \Phi(A)=\underline{G(1)}[1]\otimes L^\vee\otimes A,
 \qquad
 \Phi^{-1}(B)=L\otimes\underline{G^\vee(-1)}[-1]\otimes B.
\]
Evaluation of $G$ against its dual gives
\[
 \Phi J(C)\simeq K(C).
\]
Define the product, associator and symmetry by transport:
\begin{equation}\label{eq:gaussianstartransport}
 A\star_\psi B=\Phi^{-1}(\Phi(A)*\Phi(B)).
\end{equation}
Pulling the constant factors through $m_!$ and applying
projection formula leaves one $\underline{G(1)}[1]$.
The remaining sheaf is
\[
 m^*L\otimes\operatorname{pr}_1^*L^\vee
          \otimes\operatorname{pr}_2^*L^\vee
 =\mathcal L_\psi((x+y)^2-x^2-y^2)
 =\mathcal L_\psi(2xy).
\]
This proves \eqref{eq:gaussianstar}.
The Artin--Schreier tensor identities use the fixed
additive Lang character; the scalar phase also satisfies
\[
 2xy+2(x+y)z=2yz+2x(y+z).
\]
Transport specifies the constant lines, shifts and symmetry
in addition to that polynomial cocycle.

The ambient unit of \eqref{eq:gaussianstartransport} is
\[
 u_\psi=\delta_0\otimes\underline{G^\vee(-1)}[-1].
\]
On the full perverse image of $J$, the normalized-constant
comparisons instead give
\[
 J(C)\star_\psi J(C')\simeq J(C\otimes C'),
 \qquad e_\psi=J(E).
\]
This image unit has full support. The transported Gysin
map gives $u_\psi\to e_\psi$; convolution with an image
object on either side is an isomorphism and supplies
its image unit constraint.
Thus the two units have a precise relationship and are
not identified with each other.

The product is exact and perverse on this image, because
it is the coefficient tensor product under these
comparisons. No global t-exactness is asserted.
Full faithfulness and exactness of $J$ are those of
\fid{F50}; the new coherent comparisons make it a
symmetric strong monoidal equivalence onto its image.

\par\smallskip\noindent\textbf{The symmetry correction.}\quad
Transported symmetry is $J$ of the usual coefficient flip.
In particular it is the identity on $e_\psi\star_\psi e_\psi$
after its unit identification.

The naive kernel swap gives a different map.
Each $J$-input is lisse shifted by $[1]$, so interchanging
two such inputs gives the Koszul sign $-1$.
After the phases combine to $L(x+y)$, the fiber involution
is $y\mapsto s-y$. It preserves normalized top compact
trace, since it is an affine automorphism.
Consequently, on the image,
\[
 c_{\mathrm{naive}}\bigm|_{J(C),J(C')}
       =-J(\operatorname{flip}_{C,C'}).
\]
For $C=C'=E$ it acts by $-1$ on the unit and cannot be its
monoidal symmetry. Transport through $\Phi$ supplies
the correction. Equivalently, $\Phi J(C)=K(C)$ has
even shift $[2]$, where ordinary convolution symmetry
has no extra sign.
Thus the kernel expression alone is not being used
to choose a commutor.

\par\smallskip\noindent\textbf{Weights and duals.}\quad
Finding~\fid{F50} establishes that $G$ has weight one.
If $C$ is pure of integer weight $w$, the lisse
coefficient of $J(C)$ has weight $w-1$ and its shift
is $[1]$, so it is perverse pure of weight $w$.
Likewise $K(C)$ has coefficient weight $w-2$ and
shift $[2]$, hence complex weight $w$
\cite[5.1.8, pp.~126--127]{BBD}.

The internal monoidal dual is $J(C^\vee)$.
Its evaluation and coevaluation are $J$ of the
coefficient maps, with target or source the image unit.
Their triangle identities follow from the proved
tensor coherence.
Literal Verdier duality on the line instead gives
\[
 \mathbb D J_\psi(C)
   =L^\vee\otimes\underline{C^\vee\otimes G}(1)[1].
\]
The Gaussian pairing
$G\otimes G_{\psi^{-1}}\simeq E(-1)$ identifies
$G_{\psi^{-1}}^\vee$ with $G(1)$, so
\begin{equation}\label{eq:gaussianthreeduals}
 J(C)^{\vee_\star}=J_\psi(C^\vee),\qquad
 \mathbb D J_\psi(C)\simeq J_{\psi^{-1}}(C^\vee),\qquad
 J_{\psi^{-1}}(C^\vee(-1))
       \simeq\mathbb D J_\psi(C)(-1).
\end{equation}
The last object is the separately specified reciprocal
system of \fid{F17}, with total eigenvalues $q/\alpha$.
Internal dual, Verdier dual and that reciprocal system
are distinct constructions. A categorical dual does not
prove the second pointwise bound.
Arithmetic Jordan blocks remain allowed.

On the image, compact cohomology is inverse to $J$.
The inverse tensor comparisons make it strong monoidal
there, compatibly with the natural Gaussian evaluation.
No global monoidality for all twisted-convolution
complexes is asserted.

\par\smallskip\noindent\textbf{The actual tensor family.}\quad
Fix a regular algebraic quantum parameter and the actual
IV two-row Kronecker category, $X=V\otimes W$ with both
factors two-dimensional.
Use IV, Definition~8.1 and equation~(61), p.~39;
Propositions~10.3 and~10.5, pp.~56--57; and
Theorem~12.1, p.~68, with its immediately following
dual-corepresentation convention \cite{IV}.
These are precisely the source hypotheses of \fid{F48}.

For the actual comodule forgetful functor $\omega$, put
\[
 \mathcal T_0(N)=J((\omega N,1)).
\]
Forgetting is exact and faithful, and the underlying
tensor product is the ordinary coefficient tensor
product. The Gaussian comparisons therefore make
$\mathcal T_0$ exact, faithful and strong monoidal,
with perverse pure weight-zero images.
For every $r$ there is a coherent identification
\[
 \mathcal T_0(X^{\otimes r})
       \simeq\mathcal T_0(X)^{\star_\psi r}.
\]
The named adjacent generators remain the original
coefficient matrices
\[
 P_i^{\mathrm{ns}}=D\Pi_i,\qquad
 D=(v+v^{-1})^2,
\]
where $\Pi_i$ is the source symmetric projector.
The generator convention of the main paper is unchanged.

Every source map is sent to $J$ of its unchanged
coefficient map. Functoriality and coherence preserve
compositions, tensor products, block inclusions,
multiplication maps and all their existing relations
as actual equalities.
Identity Frobenius also makes every specified ordinary
left operator a Weil morphism, even when it is not
a comodule intertwiner.
Thus the named left and right operators of \fid{F43}
retain their matrices, commutations and original
row/opposite-composition convention as actual perverse
morphisms. These extra maps remain decorations;
they have not been declared source comodule maps.

In IV's Hopf case, existing rigid dual maps and
determinant inversion are retained using the internal
dual above. In VII's general bialgebra case, no antipode
is supplied when the determinant can vanish.
Specialization retains the actual defined action,
including any kernel it already has; faithfulness
on source morphisms does not repair a specialized
tensor action that kills the radical.
The source need not be symmetric.
A quantum Hecke braid or designated projector remains
its actual named morphism, without being identified
with the target's symmetry or exterior-power operation.

The finite-order character functor $\mathcal R_\chi$
of \fid{F48} similarly gives
$\mathcal T_\chi=J\circ\mathcal R_\chi$, exact,
faithful, strong monoidal and perverse pure of weight zero.
For example its covariance
$FE_VF^{-1}=\zeta^2E_V$ becomes the actual map
\[
 J(L_{\zeta^2})\star_\psi\mathcal T_\chi(X)
       \longrightarrow\mathcal T_\chi(X).
\]
The other raising and lowering operators have the
specified inverse phases.
Both choices keep the quantum parameter fixed;
neither imposes extension-wise parameter retuning.

\par\smallskip\noindent\textbf{The original path sign.}\quad
Nothing changes the objects $J(C)$ or their ordinary
external products. If $\gamma$ is Frobenius on $G$,
\[
 t_{C,n}(x)=-\gamma^{-n}\Tr(F^n\mid C)\psi_n(x^2),
 \qquad \sum_x t_{C,n}(x)=\Tr(F^n\mid C).
\]
For every good even path length of \fid{F50},
\[
\begin{aligned}
 R\Gamma_c(H_k,J(C)^{\boxtimes k}|_{H_k})
 \simeq{}& C^{\otimes k}\otimes G^{\vee\otimes k}\\
 &\otimes
 R\Gamma_c(H_k,\mathcal L_\psi(\textstyle\sum_i x_i^2))[k].
\end{aligned}
\]
The unshifted geometric factor has only degree $k-1$,
so the result has only degree $-1$.
Its twisted trace is exactly
\[
 q^nN_{k,C,n}
   =-\Tr(F^n\mid
       C^{\otimes k}\otimes G^{\vee\otimes k}\otimes B_{k,\psi}),
\]
with the effective space of \fid{F50}, zero error,
and an unbounded good even set.
A coefficient map $f$ induces $f^{\otimes k}$ on
the displayed coefficient factor. This proves the
claimed naturality in actual named maps and equations.
The coherent convolution comparison involves the
same coefficient tensor power.
The path itself is still the ordinary external-product
restriction required in \fid{F01}--\fid{F04};
it has not been replaced by a convolution path.
\end{proof}

\paragraph{What the operation does not change.}
The sign statement is about identity traces of effective
objects. An arbitrary inserted trace $\Tr(TF^n)$ need
not have effective spectral multiplicities, and its
tensor-inserted path output need not have the required
sign. Retaining $T$ and its tensor powers as morphisms
does not remove that distinction.

For each nonzero coefficient object, the exact pointwise
thresholds are still $b_{\max}-1$ and $1-b_{\min}$.
For one integer target $w$, the required two bounds
hold under every embedding precisely when
\[
 b_{\max,\sigma}\le w\le b_{\min,\sigma}
                    \quad\text{for every }\sigma.
\]
Thus they still assert the original pure weight.
For \fid{F45}'s total length-$L$ word object the
dual gap remains exactly $4L$.
The identity and finite-order families above already
have weight zero; their successful geometric operations
do not supply an unknown coefficient estimate or the
designated quantum-grading comparison.

The source functor is not full. For identity Frobenius,
\[
 \End(\mathcal T_0(X))\simeq M_4(E),\qquad
 \End_{\mathrm{source}}(X)=E.
\]
Its exact Grothendieck-ring character is
\[
 [\mathcal T_0(N)]=\dim_E(\omega N)[e_\psi],
 \qquad [e_\psi]\star[e_\psi]=[e_\psi].
\]
It is just the dimension character.
Geometrically every $J(C)$ is a sum of copies of the
single Gaussian lisse $[1]$ constituent.
A source simple $S$ therefore contributes $\dim S$
copies of that one constituent, rather than a new
distinguished geometric simple label.
Tensor faithfulness does not identify the source
simple basis with a KL-style perverse basis.

For \fid{F43}'s graded coefficient kernels, entrywise
$J$ and $\star_\psi$ transport direct sums, tensor
factors, and every already specified coefficient
map and word composition.
Odd degrees have odd integer Weil weights; integral
Tate motives are not asserted for them.
This does not turn a Laurent-matrix or graded-$K_0$
equality into a uniquely chosen natural isomorphism,
or provide coherence for relations established only
at $K_0$ level. The actual \fid{F48} matrix relations
are existing equalities of linear maps and are
preserved exactly.
No prescribed integral lattice, bar globalization,
cell or canonical-basis identification follows.

% Separate paragraph for Part III.
\paragraph{\fid{F53} in Grothendieck's language.}
The scheme is the additive line with its group law.
Convolution is compact pushforward along addition;
the normalized constant object is its dualizing
complex $E(1)[2]$ with compact trace $E$.
The origin Gysin map compares the ambient delta unit
with its idempotent image unit.
These are the actual Gysin, compact-trace, Kunneth
and duality constructions at the exact Milne locators
given in the proof, within the SGA compact-support
and trace frameworks used in \fid{F50}.

The Gaussian gauge tensors with a specified
Artin--Schreier line and an invertible constant
complex. The resulting kernel is the sheaf for
$2xy$, with an explicitly transported associator
and the symmetry correction required by its odd
shifts. No unspecified nonstandard sheaf theory
is needed for this operation.

The source remains the actual bialgebra/comodule
family. Its module/sheaf language is EGA~I,
Theorem~1.3.7, Corollary~1.3.8 and Theorem~1.4.1;
its coaction language is SGA~3, I, 4.7.2--3,
with cote~90, scan~17 as the audited archive anchor
\cite{EGA,SGA3,Archive}.
Those are framework references, not sources for
this particular convolution deduction.
The construction supplies coherent operations and
actual maps through all degrees together with the
original opposite-sign path cohomology.
The designated grading, full-source-Hom, simple-label,
canonical-basis and independent pointwise-weight
requirements remain distinct.

% Bibliography enrichment for the existing Milne key:
% add Theorem16.1 p108; Theorem16.7 and preceding adjunction p111;
% Theorem24.1(a),(b) p144. Other !, Kunneth and projection-formula
% locators were already supplied with F50.

